\scanpage{183}
\mspage{(I bis)}
\hfill \unsure{Mai} 1986\par
\hfill Juillet 1986

$\unsure{A|\mathcal A}$ alg. d'Az. de rang 4 sur $k$, on veut expliciter sa structure multiplicative en termes de la forme $\det$, de l'élément $1$, et de la structure \uline{spéciale} orthogonale $e\in\bigwedge^4 A$ \deleted{associée} \inserted{can.} associée à $\det$ ($e\otimes e=\mathrm{dis}(\det)$). Comme on connait la forme $\operatorname{Tr}xy$ par la formule
% slanted
\margin{\uline{NB} $\operatorname{Tr}$ \unsure{connue|connu} en termes de $\det$ et de $1$, comme $\operatorname{Tr}(x)=\psi_{\det}(x,1)$}
\[
 \det(x\unsure{+|-}y)=\det x+\det y+\operatorname{Tr}x\operatorname{Tr}y-\operatorname{Tr}xy
\]
i.e.
\[
 \psi_{\det\deleted{\unsure{(x,y)}}}(x,y)=\operatorname{Tr}x\operatorname{Tr}y-\underbrace{\operatorname{Tr}xy}_{\varphi(x,y)},
\]
\deleted{pour} et que celle-ci est non dég., pour connaître $xy$, il suffit de connaître
\[
 \operatorname{Tr}xyz\qquad\forall z\in A.
\]

Donc je cherche une formule pour $\operatorname{Tr}xyz$. Comme $\operatorname{Tr}(xy)z$ et $\operatorname{Tr}(yx)z$ ne sont pas égaux en général, i.e la structure mult. \inserted{de $A$} n'est \uline{pas} commutative
\deleted{\unsure{(Ainsi|(alors}} On sent que $\det,1$ sans $e$ ne suffisent à expliciter $\operatorname{Tr}xyz$ car ils ne déterminent la structure multiplicative que mod passage à l'inverse, i.e. mod l'involution
\[
 x\longmapsto x'=\operatorname{Tr}(x).1-x.
\]

On note aussi que $\operatorname{Tr}x\unsure{,}yz$ ne dépend que du triple $\{x,y,z\}$ et d'un ordre circulaire sur \unsure{l'ens|l'un|l'ens.} des trois termes. On a, \deleted{pour} si
\[
\begin{aligned}
\tau_{\unsure{,}}\deleted{\unsure{=}}&=\operatorname{Tr}xyz\quad(=\operatorname{Tr}yzx=\operatorname{Tr}zxy)\\
\tau'&=\operatorname{Tr}(zyx)(=\operatorname{Tr}(yxz)=\operatorname{Tr}(xzy))
\end{aligned}
\]
% archive sheet number: 182
